% Lösungen Einheit 2 von 4 – Abstand Punkt–Ebene (zusammenbau v0.1)
\einheitenkopf{Einheit 2 von 4 · Abstand Punkt–Ebene}

\erg{12}{a) Abstand berechnen \quad b) $|\vec n| = 3$, $d = \tfrac{|8 - 0 + 6 - 5|}{3} = 3$ \quad c) Lotgerade: k: $\vec x = (2 | 1 | 1) + t \cdot (2 | 1 | 4)$ \quad d) $|t| \cdot |(2 | -2 | 1)| = 3|t| = 6$, $t = 2$ oder $t = -2$: (5 | −4 | 4) oder (−3 | 4 | 0) \quad e) X(r) = (2 | 3 | 6r), $0 \le r \le 1$; $\tfrac{|12r - 12|}{3} = 2$ ergibt $r = 0{,}5$ oder $r = 1{,}5$ (nicht auf der Strecke), also X(2 | 3 | 3) \quad f) „Für welche Werte von k hat der Punkt P(1 | 3 | 1) von der Ebene $E_k$ den Abstand 1?“ – im Zähler ist P in die Ebenengleichung eingesetzt, im Nenner steht der Betrag des Normalenvektors (k | 2 | −1)}
\erg{13}{a) Fehler: nicht durch den Betrag $|\vec n| = 3$ des Normalenvektors geteilt. Richtig: $d = \tfrac{9}{3} = 3$ \quad b) Nach dem Teilen ist der Normalenvektor ein Einheitsvektor; der eingesetzte Wert ist dann die Länge der Projektion eines Verbindungsvektors zur Ebene auf die Normale, also die Länge des Lots. \quad c) $d = \tfrac{|12 + 36 - 40|}{5} = 1{,}6$ m $< 2$ m – nein}
